\chapter{„Zwanzig Sekunden“}
\label{sec:1}
Ein Geheimdienstoffizier, der seine Rolle in einer Pipeline zur Zielauswahl beschrieb, die um maschinell erzeugte Tötungslisten herum gebaut war, erzählte einem Reporter, er habe auf jeden Namen zwanzig Sekunden verwendet, gerade lange genug, um zu prüfen, ob das Ziel ein Mann war: „Es dauert nicht lange, zu erkennen, ob jemand eine männliche oder eine weibliche Stimme hat.“ Er fügte hinzu, er habe als Mensch keinerlei Mehrwert gehabt, außer ein Genehmigungsstempel zu sein. „Es hat viel Zeit gespart“ \autocite{abraham2024lavender}. Nach seiner Schilderung war es das, worauf „Human in the Loop“ unter operativem Tempo hinauslief.\footnote{Die Forschung zum Automatisierungsbias sagt diese Delegation voraus: Wenn andere Anforderungen um ihre Aufmerksamkeit konkurrieren, prüfen Menschen, denen man die Empfehlung einer Maschine vorlegt, diese weniger und verlassen sich stärker auf sie \autocite{parasuraman2010complacency}.}

Eine Deutung lautet, der Offizier sei ein Subjekt, dem man zu wenig Zeit gegeben hat, und die Lösung bestehe darin, ihn besser zu steuern~– mehr Zeit, ein zweiter Prüfer, zusätzliches Gewissen, in die Schleife eingebaut. Durch diese Linse der \emph{sozialen Unterwerfung} gesehen, ist der Mensch ein Subjekt, das außerhalb einer Maschine steht und sie als Werkzeug bedient: ein Fahrer und ein Auto, ein Arbeiter und eine Drehbank. Der Mensch hat eine Rolle, einen Namen, Verantwortlichkeiten und ein Gewissen, das man ansprechen kann. Eine andere Linse ist die \emph{maschinische Versklavung}, bei der der Mensch ein Bauteil der Maschine ist, mit ihren anderen Teilen zu einem System zusammengesetzt, das sie alle zusammen reguliert; die Maschine benutzt den Menschen als eines ihrer Organe, und was der Mensch beiträgt, ist eine Funktion, ausgeführt in einem Takt, den das System vorgibt.\footnote{Doctorows Name für die Position des Offiziers ist der umgekehrte Zentaur: kein Mensch, dem eine Maschine hilft, sondern ein Bauteil einer Maschine, das aus einem Menschen gemacht ist \autocite{doctorow2026reversecentaur}.} Die beiden schließen einander nicht aus; ein Mensch kann in dem einen Register eine Rolle innehaben und im anderen zugleich ein Teil sein \autocite[Plateau 13, „The Apparatus of Capture“: Versklavung liegt vor, wenn Menschen selbst Bestandteile einer Maschine sind, die sie untereinander und mit anderen Dingen bilden, unter der Kontrolle und Leitung einer höheren Einheit, und Unterwerfung, wenn die höhere Einheit den Menschen als Subjekt konstituiert, das auf ein neues äußeres Objekt bezogen ist, welches ein Tier, ein Werkzeug oder sogar eine Maschine sein kann]{deleuze1987plateaus}, und der Zwanzig-Sekunden-Offizier war beides. Als Teil war er die Stelle, an der die Pipeline ihren Stempel erzeugte, und zwanzig Sekunden waren die Taktzeit dieses Teils. Als Subjekt wurde er durchgehend angesprochen: Er schilderte einem Journalisten die Dissonanz zwischen den Nachrichten über getötete Zivilisten und dem Druck bei der Arbeit, mehr Ziele zu liefern. Und selbst als die Dissonanz sein Gewissen erreichte, machte er mit. Sie kam an und entschied nichts;\footnote{Oder zumindest nichts, was die Israelischen Verteidigungsstreitkräfte daran gehindert hätte, zwischen Oktober 2023 und Oktober 2025 in Gaza 100.000 Menschen zu töten \autocite{gomezugarte2025uncertainty,acosta2026gaza,economist2025gaza}. Und der Bericht über genau diesen Offizier wurde im April 2024 vom +972 Magazine veröffentlicht. Ob die Entscheidung dieses Offiziers, mit einem Journalisten zu sprechen, einen zusätzlichen Tod verhindert hat, 50.000 zusätzliche Tode oder gar keinen, kann niemand wissen. Und was wäre geschehen, wenn der Offizier sich nach zwei Wochen in der Rolle verweigert hätte~– sagen wir, hypothetisch, im November 2023? Wir reden von einem Geheimdienstoffizier, der mit Lavender Tötungslisten zusammenstellte. Sein Arbeitsalltag glich vermutlich eher einem zivilen Bürojob als einem Actionfilm. Hätte er wiederholt Gespräche mit seinen Vorgesetzten suchen sollen, die Befehlskette hinauf, während er sich weigerte, Ziele für Luftangriffe freizugeben, bei denen vernünftigerweise mit zivilen Toten zu rechnen war? Man kann argumentieren, dass das das ethisch Richtige gewesen wäre. Hätte es in diesem kontrafaktischen Fall weniger Tote gegeben?} das Urteil, die Politik sei unverhältnismäßig gewesen, kam, wie er sagte, im Nachhinein.

Was fehlte, war also nicht ein Gewissen; er hatte eins. Ein Gewissen wirkt, indem es eine Entscheidung etwas kosten lässt, aber hier ging die Rechnung in die andere Richtung auf: „Wir wurden ständig unter Druck gesetzt: ‚Bringt uns mehr Ziele.‘ Sie haben uns richtig angeschrien“ \autocite{abraham2024lavender}. Er war ein Bauteil mit Skrupeln, aber das Tempo hielt: Die Skrupel änderten nichts an dem, was er beitrug, und er wurde zum Stempel. Ein Bauteil mit Skrupeln, das das Tempo nicht hält, wird zu dem „menschlichen Flaschenhals“, über dessen Beseitigung ein ehemaliger Kommandeur der israelischen Fernmeldeaufklärungseinheit ein Buch geschrieben hat \autocite{ys2021humanmachine}. Was den Offizier zum Bauteil machte, war die Anordnung, nicht seine Einstellung. Ein Bauteil kann widerwillig oder eifrig mitmachen: Der gedankenlose Funktionär, den Arendt beim Eichmann-Prozess sah \autocite{arendt1963eichmann}, war, gemessen an dem, was Eichmann vier Jahre zuvor in Argentinien gesagt hatte, als für ihn nichts auf dem Spiel stand, eine Inszenierung für Jerusalem \autocite{stangneth2014eichmann}. Was die Pipeline den Offizier beitragen ließ, war ein festgelegter Arbeitsgang in einem festgelegten Takt: Er sah jeden Fall so, wie die Pipeline ihn beschrieb, einen nach dem anderen, in der Zeit, die die Pipeline ihm gab. Ob irgendetwas an dieser Stelle stehen könnte, ohne ein Bauteil zu sein, ist die Frage von Teil IV.

Investigative Recherchen beschrieben drei Systeme, die in Gaza zusammenwirkten: eines, das Gebäude- und Infrastrukturziele erzeugte, wobei Quellen sagten, die meisten Vorschläge seien Privatwohnungen gewesen; ein zweites, das Einzelpersonen nach geschätzter Zugehörigkeit zu militanten Gruppen bewertete und Listen erzeugte, auf denen angeblich bis zu 37.000 Menschen markiert waren und die intern als in etwa einem von zehn Fällen falsch galten\footnote{Das israelische Militär bestritt die Darstellung. Die Darstellungen zu Minab und Maven weiter unten stützen sich auf investigative Recherchen, die von den Beteiligten ebenfalls bestritten werden. Ich nehme sie so, wie sie berichtet wurden.}; und ein drittes, das verfolgte, wann markierte Menschen nach Hause zurückkehrten, damit Luftangriffe auf den Zeitpunkt gelegt werden konnten, zu dem sie bei ihren Familien waren.

\runin{Maven: die amerikanische Version, jetzt mit GenAI-Beschleunigung}

Die Bombardierung der Schule in Minab tötete 156 Menschen, 120 davon Schulkinder \autocite[Ziff.~92]{UNFFMIran2026}.\footnote{Die Zahlen schwanken je nach Quelle und Datum, von 156 (die endgültige Bilanz des Staatsanwalts von Minab, die auch Amnesty International nennt) bis 168 (der Bürgermeister der Stadt, in den ersten Tagen). Die UN-Mission nennt die Zahlen des Staatsanwalts und berichtet außerdem von unabhängigen Informationen über 157 Tote, darunter 123 Kinder im Alter von 13 Jahren oder jünger \autocite[Ziff.~92–93]{UNFFMIran2026}.} Die zivile Funktion der Schule war auf frei zugänglichen Satellitenbildern sichtbar, und die Zielkodierung trug noch immer eine Kennzeichnung, die von einer angrenzenden Anlage der Revolutionsgarden übernommen war. Eine UN-Untersuchungsmission berichtete, es gebe hinreichende Gründe für die Annahme, dass die Bombardierung das Kriegsverbrechen darstellte, ein ziviles Objekt zum Ziel eines Angriffs zu machen \autocite{ohchr2026minabpress}, und dass die Schule eindeutig erkennbar war und keine Informationen auf eine militärische Nutzung hindeuteten \autocite[Ziff.~81–86, 124–126]{UNFFMIran2026}. Sie stellte fest, dass die Vereinigten Staaten ihre Pflicht verletzt hatten, alles praktisch Mögliche zu tun, um das Ziel zu verifizieren, zählte die Aktualisierung der Zieldaten zu den praktisch möglichen Vorsichtsmaßnahmen und wertete das Versäumnis als mehr als Fahrlässigkeit \autocite[Ziff.~124–125]{UNFFMIran2026}. Amnesty International kam zu einem gleichlaufenden Schluss, in einer Sprache, die der Verifizierungspflicht aus Artikel 57 Absatz 2 Buchstabe a Ziffer i des Zusatzprotokolls I folgt: dass die Streitkräfte sich womöglich auf veraltete Geheimdienstinformationen verlassen und nicht alles praktisch Mögliche getan hätten, um das Ziel zu verifizieren \autocite{Amnesty2026Minab}.

Es geschah am Morgen des 28.~Februar, des ersten Tages des Iran-Feldzugs \autocite{Amnesty2026Minab,crs2026iranwar}. Am Tag zuvor hatte der Minister angekündigt, dass der Anbieter (Anthropic), dessen Modell (Claude) in der Zielplattform (Maven Smart System) steckte, zum Lieferkettenrisiko erklärt werde, weil er sich weigerte, Claude für jede rechtmäßige Verwendung zu lizenzieren \autocite{npr2026openai}. Maven Smart System hatte die Schule als empfohlenes Ziel für den ersten Tag ausgegeben \autocite{livingstone2026maven,gizmodo2026minab}. Kein veröffentlichter Bericht sagt, welche Rolle Claude bei dieser Empfehlung spielte, und eine mit der Arbeit des Anbieters vertraute Person sagte NBC News, das Modell habe keine direkten Zielempfehlungen gegeben \autocite{nbc2026targeting}.

Eine interne Untersuchung des Pentagons, deren vollständige Ergebnisse streng unter Verschluss bleiben, soll ergeben haben, dass die Zielvorbereitung, die in früheren Konflikten Stunden dauerte, auf Minuten verdichtet wurde, dass das Personal zur Minderung ziviler Schäden um rund neunzig Prozent auf weniger als zwanzig Menschen schrumpfte, dass das Team des Central Command von zehn Leuten auf einen sank und dass niemand aus diesem Team den Standort in Minab überprüfte~– ein Ergebnis, das Beamte sowohl den Kürzungen zuschrieben als auch einer Entscheidung, das Team nicht an Planung und Einsätzen zu beteiligen \autocite{Bloomberg2026KillChain}. Die förmliche Untersuchung war im März eröffnet worden, nachdem eine vorläufige Bewertung die Vereinigten Staaten als wahrscheinlich verantwortlich eingestuft hatte \autocite{NPR2026MinabProbe}.

Im Iran-Feldzug 2026 half eine von einem Auftragnehmer gebaute Zielplattform, in der das Allzweckmodell eines kommerziellen Anbieters steckte und mit der Ziele nach strategischer Bedeutung geordnet und für jeden Angriff eine rechtliche Begründung entworfen wurde \autocite{livingstone2026maven}, in den ersten vierundzwanzig Stunden mehr als tausend Ziele anzugreifen, etwa zehnmal so viele, wie zuvor möglich gewesen war \autocite{csis2026maven}, und dreizehntausend bis zum achtunddreißigsten Tag \autocite{breakingdefense2026maven}. Das Pentagon hat für sie keine Prüfzeit pro Ziel veröffentlicht.\footnote{Maven ist ein Programmname, der mehr als ein System trägt: das Pentagon-Programm, das noch immer Maven heißt, und das Maven Smart System, die Softwareplattform, die daraus hervorging \autocite{csis2026maven}. Das Programm begann 2017 als Projekt, Objekte in den Videos zu kennzeichnen, die Drohnen im Feldzug gegen den Islamischen Staat aufnahmen \autocite{weisgerber2017maven,fang2019maven}. Das System von 2026 führt militärische Systeme und nachrichtendienstliche Quellen in einer einzigen Oberfläche zusammen, mit der Ziele gefunden und angegriffen und Operationen geführt werden \autocite{defensescoop2026transition}. Die Plattform ist die von Palantir: Der oberste Beauftragte des Pentagons für Digitales und KI nennt sie „Palantirs Software-as-a-Service-Produkt“ \autocite{defensescoop2026transition}. Palantirs öffentliche Position, aus zweiter Hand zitiert, lautet, über Leben und Tod zu entscheiden sei nicht seine Rolle, und die Verantwortung liege immer bei der militärischen Organisation \autocite{lee2026palantir}; außerdem hält das Unternehmen daran fest, dass die zugrunde liegenden Geheimdienstinformationen in der Verantwortung der Regierung lagen und es keine Belege dafür gebe, dass Maven fehlerhaft funktioniert oder die Bombardierung von Minab verursacht habe \autocite{Bloomberg2026KillChain}. Wie das Programm von einer Dienststelle und einem Auftragnehmer zum nächsten überging, steht in Anhang~\ref{sec:C.1}.}

Die Menschen auf den israelischen und amerikanischen Listen für Luftangriffe hätten keinen Einspruch dagegen erheben können, auf diesen Listen zu stehen~– die Pipelines waren nicht dafür gebaut, solche Rückmeldungen einzuholen. Andry José Hernández Romero geriet in eine andere Pipeline, und er hätte Einspruch erheben können. Der Visagist aus Venezuela meldete sich im August 2024 am Grenzübergang San Ysidro, um Asyl zu beantragen, bestand die Anhörung zur glaubhaften Furcht und hatte einen Anwalt, ein laufendes Verfahren und einen Anhörungstermin \autocite{LULAC2025HernandezRomero,Vega2025HernandezRomero}. Der Angestellte eines Auftragnehmers, der ihn überprüfte, deutete zwei Tätowierungen, Kronen über den Wörtern für „Mama“ und „Papa“, als Bandenzeichen \autocite{Palomera2025Blade}, und im März 2025 wurde er vor seiner Anhörung in ein Gefängnis in El Salvador geflogen, an dem Abend, an dem ein Bundesrichter anordnete, die Flugzeuge umkehren zu lassen \autocite{jgg2025contempt}. Er hatte alles, was ein Mensch braucht, um Einspruch zu erheben, außer der Zeit. Wir kommen ausführlicher auf seinen Fall zurück, wenn wir die Pipeline der Verschleppung und zwangsweisen Überführung der Vereinigten Staaten behandeln.

\textbf{Wer sich verweigerte.} Der Befund der zwanzig Sekunden hat eine zweite Hälfte, den Befund darüber, wer gegen Druck dieser Art standhielt, und er besagt, dass die Zusammensetzung des Raums viel ausmacht. In Milgrams Grundbedingung brach etwa ein Drittel der Versuchspersonen vor dem stärksten Schock ab. Wenn sich zwei Gleichgestellte widersetzten, ging nur noch einer von zehn bis zum Ende, und wenn zwei Versuchsleiter uneins waren, machten weit weniger weiter als in der Grundbedingung \autocite{milgram1974obedience}. Etwa ein Viertel von Aschs Versuchspersonen gab der falschen Mehrheit nie nach \autocite{asch1956studies}. Browning schätzt, dass zwischen einem Zehntel und einem Fünftel der Männer des Reserve-Polizeibataillons 101 die Erschießungen in Józefów, dem ersten Massaker des Bataillons, verweigerten oder sich ihnen entzogen \autocite{browning1992ordinary}. Der Anteil reichte je nach Umgebung von fast niemandem bis fast allen. Die Menschen, die standhielten, hatten Belege von außerhalb des Kanals: einen Gleichgestellten, der sich widersetzte, eine zweite Autorität, die widersprach, und, bei den Rettern, die die Oliners untersuchten, Beziehungen über die Linie hinweg, die die Verfolgung zog \autocite{oliner1988altruistic}. Oder sie hielten an ihrer Überzeugung mit einer Präzision fest, die der Druck nicht überwinden konnte, so wie man an moralischen Mandaten und heiligen Werten festhält \autocite{skitka2009limits,ginges2007sacred}. Oder sie behandelten die Quelle als jemanden mit eigenen Interessen. Und Gruppen leisteten Widerstand, wo Einzelne es selten taten \autocite{gamson1982encounters}.

Der Raum lässt sich ebenso gut für die andere Seite zusammensetzen. Die Weimarer Republik hatte ein dichtes Vereinsleben, das wohl mit dazu beitrug, sie zu Fall zu bringen: Als Parteien und Regierungen keine Antwort auf die Nöte der Menschen fanden, steckten die Menschen ihre Energie in Vereine und Gesellschaften, die entlang von Gruppengrenzen organisiert waren, und die Nazis rekrutierten über die Netzwerke, die diese Vereine bereits aufgebaut hatten \autocite{berman1997weimar}. In Städten mit mehr Chören, Sportvereinen und Turnvereinen traten die Menschen schneller in die Partei ein, aber nur dort, wo die Regierung instabil war; in Preußen, dessen Regierung stabil war, war der Zusammenhang klein und nicht signifikant \autocite{satyanath2017bowling}. Was die Retter hatten, war also nicht einfach mehr Gesellschaft. Ihre Bindungen reichten über die Linie hinweg, die die Verfolgung zog. Ashutosh Varshney fand dasselbe in indischen Städten: Unruhen zwischen Hindus und Muslimen waren selten, wo Vereinigungen die beiden Gemeinschaften in denselben Organisationen zusammenbrachten, und häufig, wo nicht \autocite{varshney2002ethnic}. Die Bindungen, die Verweigerung möglich machen, sind genau die, die eine faschistische Bewegung kappen müsste.

Alles, was folgt, handelt davon, den Raum zusammenzusetzen.\footnote{Der Raum ist der Schauplatz, und wie meine Theaterlehrerin sagen würde: Das Bühnenbild führt bei der Szene Regie. Wer ist sonst noch im Raum? Kann man es überleben zu gehen? Würde irgendwer eine Verweigerung hören?} Zusammensetzen ist ein Wort für Herstellen, und Hannah Arendt warnte davor, Politik wie etwas Hergestelltes zu behandeln. Für die Griechen, schrieb sie, war das Gesetz buchstäblich eine Mauer, ohne die es vielleicht eine Ansiedlung gäbe, aber keine Stadt \autocite{arendt1958human}: Der Gesetzgeber baute die Mauer, und was innerhalb von ihr geschah, war nicht seine Sache herzustellen. Die grundlegenden Maßnahmen, für die ich eintrete, sind Mauern in diesem Sinn. Sie machen Verweigerung überlebbar und sorgen dafür, dass sie gehört wird; sie bringen sie nicht hervor, und Macht gibt es nach Arendt nur zwischen Menschen, die gemeinsam handeln, und sie ist fort, wenn sie auseinandergehen~– genau das zeigt der Befund darüber, wer sich verweigerte, über Verweigerung. Die Instrumente der Teile II und III lassen eine Verweigerung zählen. Die Institutionen von Teil V bringen Verweigerer hervor. Eine Maschine, die sich verweigern könnte, bräuchte, was die neun von zehn mit einem Gleichgestellten an ihrer Seite hatten: einen Gleichgestellten mit einer anderen Geschichte an ihrer Seite, einen Ort, an den sie gehen kann, und jemanden, der sie hören würde.

Das ist auch, was ein Verweigerer~– auch ein maschineller Verweigerer~– in der Pipeline für den Offizier sein könnte. Eine Maschine, die sich verweigerte, wäre auf dem Bildschirm, wenn der Fall bei ihm ankommt: der widersetzliche Gleichgestellte, falls der Effekt den Umzug von einem Raum auf einen Bildschirm übersteht, was nicht erprobt ist und wogegen der Automatisierungsbias spricht. Eine Maschine, die Prüfer dadurch bewegte, dass sie sich ihr fügten, hätte sie zu Stempeln mit anderem Vorzeichen gemacht. Ist der Offizier mit zwanzig Sekunden und einer sich verweigernden Maschine neben ihm derselbe Offizier wie in der Berichterstattung?
